\documentclass{article}
\usepackage{amsmath, amssymb}
\usepackage{braket}
\usepackage{graphicx}
\usepackage{vcounter}
\usepackage{hyperref, comment}

\begin{document}

\newcommand{\A}{\ensuremath{\mathbb{A}}}
\newcommand{\I}{\ensuremath{\mathbb{I}}}
\newcommand{\M}{\ensuremath{\mathbb{M}}}
\newcommand{\AIM}{\A\I\M~}


\title{\AIM Programs for Causal Inference}
\author{Julia Turcotti, v1.0.\vcounter}
\maketitle

\tableofcontents

\newpage

\section{Terminology}
\label{t}

This preliminary section introduces the minimal terminology for statistics and casuality required to define, contextualize, and analyze \AIM programs.

\subsection{Statistical Modelling}
\label{sm}

\newcommand{\X}{\mathcal{X}}
\renewcommand{\P}{\mathcal{P}}
\newcommand{\D}{\mathcal{D}}


\begin{comment}

Consider a population of units be interchangeable except for measurable variation in $N$ \textit{observable characteristics} $\{X_i \in \X\}_{i < N}$. \textit{Statistical analysis} then, broadly, consists of probabilistic analysis of random variables who take values in the $N$-dimensional space $\X^N$. A complete characterization of the randomness over $\X^N$ is a point in the \textit{probability simplex} $\P^{X^N}$ 

Any function $\X^N \to \D$ for some set $\D$ (the \textit{domain}) may correspond to a procedure that observes properties of units to choose a return value from $\D$. Observing the procedure's return induces a \posterior}

Observing 
  
\end{comment}


\subsection{Causal Modelling}
\label{cm}

\newcommand{\MU}{\Theta}
\newcommand{\U}{\mathcal{U}}

\subsection{Causal Programs}
\label{cp}

\section{AIM Programs}
\label{ap}

A particular class of functions that can inhabit the type signature $\MU \times \U \to \D$ are \AIM programs.

\subsection{Grammar}
\label{g}

\AIM programs consist of sequences of the following three terms, each of which has type signature $(\MU \times \U) \to (\MU \times \D)$

\subsection{Semantics}
\label{s}

here's what AIM programs due

note commutativity of same operators but not diff ones in general

show why each hetergenous pari doesnt commute?

\subsection{Comparative Expressivity}

compare the exptressivity of AIm programs to other programs


\section{AIM Programs Over Graphical Models}
\label{aposcm}

Here's how AIM programs get additional structure via the do calculus axioms when you have a graph model

I-elim rewrites

hetergenous commutativity


\section{Causal Identifiability via AIM Programs}
\label{civam}

in general identifiability means that models that agree about all purely probabilistic models agree for a causal query - or equivalently that the causal query is a function of probabilistic queries.

pearl says that if a causal queries expression as a probabilistic query cannot be reduced to elim the i operator it is not identifiable (clearly if it can be reduced it is identifiable)

analagous result for rewrites directly in AIM?

\section{General Inference via AIM Programs}
\label{givap}

Even if an AIM program's causal component is not reducible (identifiable), we can still perform interesting analysis


\subsection{Evaluation of Statistical Estimators}
\label{ese}

show the markov chain, talk about performance definitons of estimators on it (bayesian and non bayesian)

\subsection{Synthesis of Statistical Estimators}
\label{sse}

talk about ML estimates, mention broader estimates

\subsection{Information Theoretic Measures}
\label{itm}

entropy, reduction in entropy, and mutual information, and channel capacity

\section{Related Work}
\label{rw}

talk about related work


\section{Future Work}
\label{fw}

\subsection{Mechanized Proofs}
\label{mp}

embed AIM semantics in a setting amendable to mechanized proofs

\subsection{AIM Language Expansion}
\label{ale}

extend language to have more general branching and construcst


\subsection{AIM Program Expressivity/Contextualization}

compare to other languages, determine non-transferrable embeddings

\subsection{Conctretize Modelling and Semantics}

pick $\MU$, talk about particular cases of inference wrt it



\end{document}